\documentclass[10pt,a4paper]{amsart}
\usepackage[margin=19mm]{geometry}
\usepackage{amsmath,amssymb,mathtools}
\usepackage[scr=rsfs,cal=euler]{mathalfa}
\usepackage{xfrac}
\usepackage[unicode,hidelinks,bookmarks=false]{hyperref}
\newcommand{\ie}{i.e.\ }
\newcommand{\cf}{cf.\ }
\newcommand{\resp}{respectively\ }
\newcommand{\Subsection}{\subsection}
\providecommand{\nobreakdash}{\nobreak\hbox{-}}
\setlength{\emergencystretch}{2em}
\multlinegap0pt
\usepackage{bm}
\usepackage{bbm}
\usepackage{dsfont}
\usepackage{xspace}
\usepackage{cancel}
\usepackage{mathtools}
\usepackage{amsthm}
\makeatletter
\@ifundefined{theorem}{\newtheorem{theorem}{Theorem}}{}
\@ifundefined{lem}{\newtheorem{lem}[theorem]{Lemma}}{}
\makeatother
\def\S{\mathfrak{S}}
\def\carre{[0,1]^2}
\def\mm{\!\!-\!\!}
\def\m{\!-\!}
\def\p{\!+\!}
\newcommand{\Binomial}[2]{{\rm Binom}\left(#1 , #2\right)}
\newcommand{\card}[1]{\left|#1\right|}
\newcommand{\perm}[1]{{\rm Perm}\left(#1\right)}
\newcommand{\prob}[1]{\mathbf{P}\left(#1\right)}
\title[A geometric approach to conjugation-invariant random permuta]{A geometric approach to conjugation-invariant random permutations}
\author{Victor Dubach}
\date{Source: arXiv:2402.10116v4 (CC BY 4.0)}
\begin{document}
\maketitle

\noindent\emph{Source and licence.} A geometric approach to conjugation-invariant random permutations. Version arXiv:2402.10116v4 (CC BY 4.0), distributed under the Creative Commons Attribution 4.0 International licence, as recorded in the arXiv licence metadata for this exact version and shown on the abstract page https://arxiv.org/abs/2402.10116v4. Full rights notice: licenses/arxiv-2402.10116-CC-BY-4.0.txt.\par\smallskip

\noindent\textit{Editorial scope.} Counted unit: the Lem whose source label is \texttt{lem: regularity psi with fixed points} in the article (source0.tex lines 2179--2182, source label \texttt{lem: regularity psi with fixed points}), delivered together with the article's own complete original proof of it (source0.tex lines 2184--2249, one \texttt{proof} environment of 66 lines). No mathematics is written, repaired or re-derived: statement and proof are the article's own text line for line. Required context retained uncounted: none. Every definition, display and label of those lines is retained unchanged. Omitted from this package and not redistributed: 1--2178, 2183--2183, 2250--2436. Nothing inside a retained excerpt is omitted or altered: every retained body is delivered exactly as the article's lines are written, so no omission receipt of any kind accompanies this package. Citations: the retained excerpts carry no citation command at all, so no bibliography entry of the article is transcribed and no reference is redistributed. Reference adaptation: none was needed and none was made; every reference inside the delivered unit resolves inside it, so no unresolved reference or citation is produced. Printed slips kept literal: no printed word, symbol, hypothesis or number of the counted statement or of its proof was altered. Layout and class: the article's own class is \texttt{article} with packages including graphicx, geometry, fontenc, times, amsmath, amsfonts, amssymb, amsthm, xcolor, hyperref, cleveref, authblk; the wrapper is the lane's pinned \texttt{amsart} editorial wrapper at 10pt on A4, which restates the theorem-like environments the retained text needs and adds the article's own macro definitions that the retained text needs (\texttt{\textbackslash Binomial}, \texttt{\textbackslash S}, \texttt{\textbackslash card}, \texttt{\textbackslash carre}, \texttt{\textbackslash m}, \texttt{\textbackslash mm}, \texttt{\textbackslash p}, \texttt{\textbackslash perm}, \texttt{\textbackslash prob}), with extra wrapper packages (bm, bbm, dsfont, xspace, cancel, mathtools). The delivered numbering is the unit's own. Assets: no retained span contains a figure environment, a graphics-inclusion command, a PDF-page inclusion, a table or a TikZ picture; no image file is copied into this package, the package declares no asset dependency and no image file is redistributed with it. Licence: version arXiv:2402.10116v4 of this article is distributed under the Creative Commons Attribution 4.0 International licence, as recorded in the arXiv licence metadata for this exact version and shown on the abstract page https://arxiv.org/abs/2402.10116v4. A full rights notice is delivered at licenses/arxiv-2402.10116-CC-BY-4.0.txt. Characters: every retained excerpt is pure ASCII, so the delivered engine, XeLaTeX with its default Latin Modern text font, reports no missing character.
\begin{lem}\label{lem: regularity psi with fixed points}
    For any $p_1\in[0,1]$ and $\pi\in\S_r$, the function $\psi_\pi^{p_1}$ is $(r\m1)$-Lipschitz with respect to the $L^1$-norm on $\carre$.
    Moreover, its restriction to $\{0\}\times[0,1]$ is a polynomial which is non-constant if $p_1<1$ and $r\ge 2$..
\end{lem}

\begin{proof}
    Write $\psi:=\psi_\pi^{p_1}$. Recall the variables $\hat Z_1,\dots,\hat Z_r$ used to define $\psi$.
    For any $z=(u,v),z'=(u',v')\in\carre$:
    \begin{equation*}
        \card{\psi(z)-\psi(z')}
        \le \prob{\perm{\hat Z_1,\dots,\hat Z_{r-1},z} \ne \perm{\hat Z_1,\dots,\hat Z_{r-1},z'}} .
    \end{equation*}
    Define $\mathcal{B}_{z,z'}$ as the set of points in $\carre$ having x-coordinate between $u$ and $u'$ or y-coordinate between $v$ and $v'$.
    For the permutations $\perm{\hat Z_1,\dots,\hat Z_{r-1},z}$ and $\perm{\hat Z_1,\dots,\hat Z_{r-1},z'}$ to differ, there needs to be a $\hat Z_i$ whose position with respect to $z$ (NW, NE, SW, or SE) is not the same as with respect to $z'$.
    Therefore:
    \begin{equation*}
        \card{\psi(z)-\psi(z')}
        \le \sum_{i=1}^{r-1}\prob{\hat Z_i \in \mathcal{B}_{z,z'}}
        \le (r\m1)\left(|u-u'|+|v-v'|\right)
    \end{equation*}
    and $\psi$ is indeed $(r\m1)$-Lipschitz.
    Now fix $v\in[0,1]$ and let $Q$ be the number of $\hat Z_i's$, $i\le r\m1$ on the diagonal $\Delta$.
    The random variable $Q$ follows a $\Binomial{r\m1}{p_1}$ law.    
    Now let $q$, $j_{\rm SW}$, $j_{\rm SE}$, $j_{\rm NE}$, $j_{\rm NW}$, $k_{\rm SW}$, $k_{\rm NE}$ be nonnegative integers satisfying
    \begin{equation}\label{eq: phi - decoupage points si diago}
        \left\{
        \begin{array}{lllll}
        q \in\{0,1,\dots,r-1\} ;\\
        j_{\rm SW}+j_{\rm SE}+j_{\rm NE}+j_{\rm NW}+k_{\rm SW}+k_{\rm NE} = r-1 ;\\
        k_{\rm SW}+k_{\rm NE} = q ;\\
        j_{\rm SW}+j_{\rm SE}+k_{\rm SW} = \pi(1)-1 ;\\
        j_{\rm NE}+j_{\rm NW}+k_{\rm NE} = r-\pi(1) ;
        \end{array}
        \right.
    \end{equation}
    (the last line is redundant, but we keep it for clarity).
    Under this condition we write $A_{q,\bf j,k}(v)$ for the event that $Q=q$; $j_{\rm SW},j_{\rm SE},j_{\rm NE},j_{\rm NW}$ are the numbers of non-diagonal $\hat Z_i$'s in each quadrant defined by $(v,v)$; and $k_{\rm SW}$, $k_{\rm NE}$ are the numbers of diagonal $\hat Z_i$'s lying southwest, resp.~northeast of $(v,v)$.
    If $\perm{\hat Z_1,\dots,\hat Z_{r-1},(0,v)}=\pi$ then $A_{q, \bf j,k}(v)$ is verified for some $q,\bf j,k$ satisfying \eqref{eq: phi - decoupage points si diago}.
    Moreover, conditionally on $A_{q,\bf j,k}(v)$, the probability of this event does not depend on $v$.
    Indeed one can construct a straightforward coupling between $(\hat Z_i)_{1\le i\le r-1}$ conditioned on $A_{q,\bf j,k}(v)$ and $(\hat Z_i')_{1\le i\le r-1}$ conditioned on $A_{q,\bf j,k}(v')$ such that a.s.~$\perm{\hat Z_1,\dots,\hat Z_{r-1},(0,v)}=\perm{\hat Z_1',\dots,\hat Z_{r-1}',(0,v')}$.
    Therefore:
    \begin{align*}
        \psi(0,v) 
        &= \sum_{{q,\bf j,k}\text{ s.t. }\eqref{eq: phi - decoupage points si diago}}
        \prob{\left.\perm{\hat Z_1,\dots,\hat Z_{r-1},(0,v)}=\pi \;\right| A_{q,\bf j,k}(v)}
        \prob{A_{q,\bf j,k}(v)}
        \\&= \sum_{{q,\bf j,k}\text{ s.t. }\eqref{eq: phi - decoupage points si diago}}
        c_{\pi, q,\bf j,k} \binom{r\m1}{q} p_1^q (1\mm p_1)^{r\m q\m 1} 
        \binom{r\m q\m 1}{j_{\rm SW},j_{\rm SE},j_{\rm NE},j_{\rm NW}}
        v^{k_{\rm SW}\p 2j_{\rm SW}\p j_{\rm SE}\p j_{\rm NW}} 
        (1\mm v)^{k_{\rm NE}\p 2j_{\rm NE}\p j_{\rm SE}\p j_{\rm NW}}
        \\&= v^{\pi(1)\m 1}(1\mm v)^{r\m \pi(1)} 
        \sum_{{q,\bf j,k}\text{ s.t. }\eqref{eq: phi - decoupage points si diago}}
        c_{\pi, q,\bf j,k} \binom{r\m1}{q} p_1^q (1\mm p_1)^{r\m q\m 1}
        \binom{r\m q\m 1}{j_{\rm SW},j_{\rm SE},j_{\rm NE},j_{\rm NW}}
        v^{j_{\rm SW}\p j_{\rm NW}} (1\mm v)^{j_{\rm NE}\p j_{\rm SE}}
        \\&= v^{\pi(1)-1}(1-v)^{r-\pi(1)} \sum_{q=0}^{r-1} \sum_{m=0}^{r-q-1}
        c_{\pi, q, m}'\, v^m (1-v)^{r-q-1-m}
    \end{align*}
    for some non-negative constants $c_{\pi, q,\bf j,k}$ and $c_{\pi, q, m}'$.
    Thus $\psi(0,v)$ is a polynomial in $v$.
    It remains to see that it is non-constant if $p_1<1$ and $r\ge2$.
    
    First suppose that $\pi(1)\ne1$.
    The polynomial $\psi(0,v)$ contains a monomial $v^{\pi(1)-1} \sum_{q=0}^{r-1} c_{\pi,q,0}'$ where $c_{\pi,q,0}'\ge0$, along with terms of higher degrees.
    Since $p_1<1$, the event $A_{0,\bf j,0}(v)$ has positive probability for any $(0,\bf j,0)$ satisfying \eqref{eq: phi - decoupage points si diago}, and also $c_{\pi,0,\bf j,0} >0$.
    In particular for $j_{\rm SW}+j_{\rm NW} = 0$, we deduce $c_{\pi,0,0}'>0$.
    Thus $\psi(0,v)$ is indeed non-constant. 
    
    On the other hand if $\pi(1)=1$ then $\pi(1)\ne r$ and we can do the same reasoning with the polynomial $\psi(0,1-v)$ to conclude.
\end{proof}

\end{document}
